% GENERATED FILE. Edit canonical lessons, then run npm run generate:books.
% canonical-source-hash: fnv1a64:83dc52d93b881643
% canonical-lessons: ES-C07-comer, ES-C07-como, ES-C07-comes, ES-C07-come, ES-C07-repaso-comer

\chapter{Comer --- the Second Verb Family, One Slot at a Time}
\label{ch:spanish-comer}
\hlchaptermodality{24}

\begin{chapteropening}
\textbf{By the end of this chapter:} \emph{I can use all three singular forms of a regular -er verb, having built each one separately.}

The -er singular, earned a slot at a time - and the yo ending that never changed.
\end{chapteropening}

\section[comer]{comer — \textquotedblleft{}to eat,\textquotedblright{} and the second verb family}

\label{lesson:ES-C07-comer}



\begin{quote}
Last chapter you learned the \textbf{-ar} verbs. Spanish has two more families — \textbf{-er} and \textbf{-ir} — and here's the good news: they're \textbf{almost the same}. Meet \textbf{comer} (\textquotedblleft{}to eat\textquotedblright{}), your first \textbf{-er} verb.
\end{quote}

\begin{sounds}
\begin{itemize}
\raggedright
  \item \texttt{r-tap} — one soft tap: \emph{ko-\textbf{MER}}.
  \item The \emph{c} before \emph{o} is a hard \emph{k}: \emph{comer} = \emph{ko-MER}.
\end{itemize}
\end{sounds}

\begin{cousinweb}
\textbf{comer} derives from Latin \textbf{comedere}, \emph{\textquotedblleft{}to eat up\textquotedblright{}} — \textbf{com-} (intensifier, \textquotedblleft{}completely\textquotedblright{}) + \textbf{edere}, \textquotedblleft{}to eat.\textquotedblright{} That \textbf{edere} is the twin of English \textbf{eat} (both from the same ancient root), and it feeds a small English family:

\begin{itemize}
\raggedright
  \item \textbf{edible} — fit to be eaten.
  \item \textbf{comestible} — a fancy word for food (straight from \emph{comedere}).
  \item English \textbf{eat} itself — the Germanic cousin of Latin \emph{edere}.
\end{itemize}
\end{cousinweb}

\begin{grammarlens}[title={Grammar Lens: one family over, and one ending already yours}]
\emph{Hablar} ended in \textbf{-ar}. \emph{Comer} ends in \textbf{-er}, and that single letter is what makes it a different family.

Here is the good news, and it is worth having before you meet a single new form: \textbf{the \emph{yo} ending does not change.} \emph{Hablar} gave \emph{habl\textbf{o}}. \emph{Comer} will give \emph{com\textbf{o}}. The ending you already own works untouched.

Only two slots in this family behave differently from the one you know, and you will meet them one at a time over the next three chapters — not here.
\end{grammarlens}

\subsection*{Your turn}
Say these aloud:

\begin{itemize}
\raggedright
  \item \textquotedblleft{}comer\textquotedblright{} — \emph{ko-MER}
  \item \textquotedblleft{}comer\textquotedblright{} then English \textquotedblleft{}edible, comestible, eat\textquotedblright{} — the \emph{edere} family
\end{itemize}

\begin{tcolorbox}[breakable,colback=teal!4,colframe=teal!35!black,title={Before you move on}]
What Latin verb is \emph{comer} from, and its English cousins? (\emph{Comedere} $\leftarrow$ \emph{edere} — edible, eat.) Which ending is the same in both families? (The \textbf{yo} ending, \textbf{-o}.) Next: the \emph{yo} form — and it will cost you nothing.
\end{tcolorbox}
\section[como]{como — \textquotedblleft{}I eat,\textquotedblright{} and you already knew how}

\label{lesson:ES-C07-como}



\begin{quote}
How do you say \textquotedblleft{}I speak\textquotedblright{}? (\emph{Hablo}.) And what is the verb for \textquotedblleft{}to eat\textquotedblright{}? (\emph{Comer}.)
\end{quote}

\begin{grammarlens}[title={Grammar Lens: the ending that did not change}]
Drop the \textbf{-er}, add \textbf{-o}:

\begin{quote}
\textbf{comer} $\to$ \textbf{como} — \emph{I eat.}
\end{quote}

That is the same operation and the same ending you used on \emph{hablar}. A new verb family, and the \emph{yo} slot cost you nothing at all.

This is worth pausing on, because it is how most of the rest will go. Spanish's three families differ in \textbf{two} of their singular slots, never in the first one. Whatever the verb, \textquotedblleft{}I\textquotedblright{} is \textbf{-o}.

And as before, one word is the whole sentence: \textbf{como}. No pronoun.
\end{grammarlens}

\begin{cousinweb}
One warning, because this word is busy. \textbf{Como} is three things:

\begin{itemize}
\raggedright
  \item \textbf{como} — \emph{I eat}, the form you just built.
  \item \textbf{como} — \emph{like}, \emph{as}. (\emph{Como tú} — \textquotedblleft{}like you.\textquotedblright{})
  \item \textbf{cómo}, with the accent — \emph{how?}, which you have used since your first conversations.
\end{itemize}

Same four letters, three jobs. The accent marks the question word; context sorts the other two. This is not a trap so much as a reminder that the accent in Spanish is doing real work.
\end{cousinweb}

\subsection*{Your turn}
Say these aloud:

\begin{itemize}
\raggedright
  \item \textquotedblleft{}hablo\textquotedblright{} then \textquotedblleft{}como\textquotedblright{} — the same ending, twice
  \item \textquotedblleft{}como\textquotedblright{} alone, as a complete sentence
\end{itemize}

\begin{tcolorbox}[breakable,colback=teal!4,colframe=teal!35!black,title={Before you move on}]
How do you say \textquotedblleft{}I eat\textquotedblright{}? (\textbf{Como}.) Which slot never changes between verb families? (The \textbf{yo} slot — always \textbf{-o}.) What does \emph{cómo} with an accent mean? (\textbf{How?}) Next: the slot that \emph{does} change.
\end{tcolorbox}
\section[comes]{comes — \textquotedblleft{}you eat,\textquotedblright{} and the first real difference}

\label{lesson:ES-C07-comes}



\begin{quote}
Say \textquotedblleft{}I eat.\textquotedblright{} (\emph{Como}.) Now say \textquotedblleft{}you speak,\textquotedblright{} to a friend. (\emph{Hablas}.)
\end{quote}

\begin{grammarlens}[title={Grammar Lens: one vowel moves}]
In the \emph{-ar} family, \emph{tú} took \textbf{-as}. Here it takes \textbf{-es}:

\begin{quote}
\textbf{comer} $\to$ \textbf{comes} — \emph{you eat.}
\end{quote}

That is the whole difference. The \textbf{-s} that means \emph{you} is still there, doing the same job it does in \emph{hablas}. Only the vowel in front of it has moved from \textbf{a} to \textbf{e} — and it moved because the verb's own family vowel is \textbf{e}. The ending is agreeing with the family it belongs to.

So you are not learning a new ending so much as watching a familiar one take on the colour of its surroundings.
\end{grammarlens}

\begin{cousinweb}
Latin had \emph{amās} for \textquotedblleft{}you love\textquotedblright{} and \emph{comedis} for \textquotedblleft{}you eat\textquotedblright{} — different vowels in the ending, exactly as Spanish has now. The split is not a Spanish invention; it is inherited, and it is why the families exist at all.

What every one of them kept, across two thousand years, is the \textbf{-s}. If a Spanish verb form ends in \emph{-s}, it is talking to \emph{tú}. That single rule will carry you through every tense in this book.
\end{cousinweb}

\subsection*{Your turn}
Say these aloud:

\begin{itemize}
\raggedright
  \item \textquotedblleft{}hablas\textquotedblright{} then \textquotedblleft{}comes\textquotedblright{} — the same \emph{-s}, a different vowel
  \item \textquotedblleft{}como\textquotedblright{} then \textquotedblleft{}comes\textquotedblright{} — and feel the sentence turn toward the other person
\end{itemize}

\begin{tcolorbox}[breakable,colback=teal!4,colframe=teal!35!black,title={Before you move on}]
What is \textquotedblleft{}you eat,\textquotedblright{} to a friend? (\textbf{Comes}.) What part of the ending is the same as in \emph{hablas}? (The \textbf{-s}.) What changed, and why? (The vowel, to match the verb's own family.) Next: the last of the three.
\end{tcolorbox}
\section[come]{come — \textquotedblleft{}he eats,\textquotedblright{} and the polite \textquotedblleft{}you eat\textquotedblright{}}

\label{lesson:ES-C07-come}



\begin{quote}
Say \textquotedblleft{}you eat,\textquotedblright{} to a friend. (\emph{Comes}.) And what was the \emph{usted} form of \emph{hablar}? (\emph{Habla}.)
\end{quote}

\begin{grammarlens}[title={Grammar Lens: take the -s off again}]
Drop the \textbf{-er}, add \textbf{-e}:

\begin{quote}
\textbf{comer} $\to$ \textbf{come} — \emph{he eats}, \emph{she eats}, and \textbf{usted come}, \emph{you eat}, politely.
\end{quote}

The relationship is the one you already know. In the \emph{-ar} family, \emph{hablas} became \emph{habla} by losing its \textbf{-s}. Here \emph{comes} becomes \emph{come} the same way.

And exactly as before, this one form does three jobs — \emph{he}, \emph{she}, and the respectful \emph{you} — for the same reason it did last time: \emph{usted} began life as \emph{vuestra merced}, a thing you talk \textbf{about}.

Three slots, three chapters, and the shape of the family is now yours.
\end{grammarlens}

\begin{cousinweb}
Latin's form here was \emph{comedit}. Spanish let the final \textbf{-t} fall, precisely as it let it fall from \emph{hablat} to give \emph{habla}.

That is the useful part: the two families were not worn down by different forces. The same weather hit both. Which is why, once you know what happened to one, the other holds no surprises.
\end{cousinweb}

\subsection*{Your turn}
Say these aloud:

\begin{itemize}
\raggedright
  \item \textquotedblleft{}habla\textquotedblright{} then \textquotedblleft{}come\textquotedblright{} — the same move, one family over
  \item \textquotedblleft{}come\textquotedblright{} for a friend's sister, then for a stranger you are being polite to
\end{itemize}

\begin{tcolorbox}[breakable,colback=teal!4,colframe=teal!35!black,title={Before you move on}]
What is \textquotedblleft{}she eats\textquotedblright{}? (\textbf{Come}.) How did it come from \emph{comes}? (By losing the \textbf{-s}, exactly as \emph{hablas} $\to$ \emph{habla}.) What did Latin's \emph{comedit} lose? (Its \textbf{-t}.) A short review closes this chapter; the third and last family follows.
\end{tcolorbox}
\section[Practice]{Review — the -er family, all three}

\label{lesson:ES-C07-repaso-comer}



\begin{quote}
Nothing new. Say the three forms of \emph{comer} you have built, in any order they come to you.
\end{quote}

\begin{grammarlens}[title={Grammar Lens: the table you have earned}]
Three slots, three lessons. Now they are worth seeing beside the family you already had:

\noindent
\begin{tabularx}{\linewidth}{@{}>{\raggedright\arraybackslash}X>{\raggedright\arraybackslash}X>{\raggedright\arraybackslash}X>{\raggedright\arraybackslash}X@{}}
\toprule
\textbf{who} & \textbf{\emph{-ar}} & \textbf{\emph{-er}} & \textbf{} \\
\midrule
I & -o & \textbf{-o} & \emph{hablo} · \textbf{como} \\
tú & -as & \textbf{-es} & \emph{hablas} · \textbf{comes} \\
él · ella · usted & -a & \textbf{-e} & \emph{habla} · \textbf{come} \\
\bottomrule
\end{tabularx}

Read the middle two columns down and the whole \emph{-er} family is one observation: \textbf{the \emph{yo} slot is shared, and the other two swap a for e.}

That is the entire cost of a second verb family. Every regular \emph{-er} verb in Spanish runs on it.
\end{grammarlens}

\subsection*{Your turn}
Say these aloud:

\begin{itemize}
\raggedright
  \item \textquotedblleft{}como, comes, come\textquotedblright{} straight through
  \item \textquotedblleft{}hablo, como\textquotedblright{} · \textquotedblleft{}hablas, comes\textquotedblright{} · \textquotedblleft{}habla, come\textquotedblright{} — the two families in step
\end{itemize}

\emph{Twice through:}

\begin{tcolorbox}[breakable,colback=teal!4,colframe=teal!35!black,title={Before you move on}]
Give the three \emph{-er} endings. (\textbf{-o, -es, -e}.) Which one matches \emph{-ar} exactly? (The \textbf{yo} ending.) How many regular \emph{-er} verbs does this fit? (\textbf{All of them}.) Next chapter: the third family, and it is the smallest change yet.
\end{tcolorbox}
